% GENERATED FILE. Edit canonical lessons, then run npm run generate:books.
% canonical-source-hash: fnv1a64:1ff1f2e6eb337f2e
% canonical-lessons: GU-C474-lal, GU-C474-cavvu, GU-C474-cusvu, GU-C474-thukvu, GU-C474-talvu

\chapter{Eating Slowly}
\label{ch:gu-eating-slowly}
\hlchaptermodality{474}

\begin{chapteropening}
\textbf{By the end of this chapter:} \emph{I can talk about chewing and swallowing.}

Complete the last lesson of chapter 474: I can talk about chewing and swallowing.
\end{chapteropening}

\section[\texorpdfstring{\gu{લાળ} (lāḷ)}{લાળ (lāḷ)}]{\gu{લાળ} (lāḷ) — saliva; drool}

\label{lesson:GU-C474-lal}



\begin{quote}
Before the new one: say the Gujarati for the calf of the leg, then the Gujarati for an armpit.
\end{quote}

\subsection*{You'll want to know: \gu{લાળ}}
\textbf{\gu{લાળ}} — \emph{lāḷ} — \textquotedblleft{}saliva; drool\textquotedblright{}.

\begin{quote}
\gu{ઊંઘમાં} \gu{બાળકના} \gu{મોઢામાંથી} \gu{લાળ} \gu{પડે} \gu{છે}.
\end{quote}

\begin{quote}
\emph{ū̃ghmā̃ bāḷaknā moḍhāmā̃thī lāḷ paḍe che.}
\end{quote}

\begin{quote}
The baby drools in its sleep.
\end{quote}

A feminine noun.

\begin{cousinweb}
\textbf{\gu{લાળ}} (\emph{lāḷ}) is inherited from Sanskrit \emph{lālā}, \textquotedblleft{}spittle, saliva\textquotedblright{}; CDIAL lists the Gujarati word (CDIAL 11027).
\end{cousinweb}

\begin{grammarlens}[title={What you've built}]
One more word for describing what you lived and felt.
\end{grammarlens}

\subsection*{Your turn}
\begin{itemize}
\raggedright
  \item \emph{Say it:} \emph{lāḷ}
  \item \emph{Say it:} \emph{lāḷ}, once more
  \item \emph{Say it:} \emph{ū̃ghmā̃ bāḷaknā moḍhāmā̃thī lāḷ paḍe che.} — The baby drools in its sleep.
  \item \emph{Recall:} say the Gujarati for the calf of the leg, then the Gujarati for an armpit, then say \emph{lāḷ} again
\end{itemize}

\begin{tcolorbox}[breakable,colback=teal!4,colframe=teal!35!black,title={Before you move on}]
What does \gu{લાળ} mean? (\textquotedblleft{}Saliva; drool\textquotedblright{}.) Say it once more.
\end{tcolorbox}
\section[\texorpdfstring{\gu{ચાવવું} (cāvvũ)}{ચાવવું (cāvvũ)}]{\gu{ચાવવું} (cāvvũ) — to chew}

\label{lesson:GU-C474-cavvu}



\begin{quote}
Before the new one: say the Gujarati for an armpit, then the Gujarati for saliva.
\end{quote}

\subsection*{You'll want to know: \gu{ચાવવું}}
\textbf{\gu{ચાવવું}} — \emph{cāvvũ} — \textquotedblleft{}to chew\textquotedblright{}.

\begin{quote}
\gu{દાદા} \gu{ધીમે} \gu{ધીમે} \gu{રોટલી} \gu{ચાવે} \gu{છે}.
\end{quote}

\begin{quote}
\emph{dādā dhīme dhīme roṭlī cāve che.}
\end{quote}

\begin{quote}
Grandfather chews his flatbread slowly.
\end{quote}

Past, with \gu{તેણે} (he or she …): \gu{તેણે} \gu{ભાત} \gu{ચાવ્યો}, \emph{teṇe bhāt cāvyo}.

\begin{cousinweb}
\textbf{\gu{ચાવવું}} (\emph{cāvvũ}) is inherited from Sanskrit \emph{carvati}, \textquotedblleft{}chews\textquotedblright{}; CDIAL lists the Gujarati verb (CDIAL 4711).
\end{cousinweb}

\begin{grammarlens}[title={What you've built}]
One more word for describing what you lived and felt.
\end{grammarlens}

\subsection*{Your turn}
\begin{itemize}
\raggedright
  \item \emph{Say it:} \emph{cāvvũ}
  \item \emph{Say it:} \emph{cāvvũ}, once more
  \item \emph{Say it:} \emph{dādā dhīme dhīme roṭlī cāve che.} — Grandfather chews his flatbread slowly.
  \item \emph{Recall:} say the Gujarati for an armpit, then the Gujarati for saliva, then say \emph{cāvvũ} again
\end{itemize}

\begin{tcolorbox}[breakable,colback=teal!4,colframe=teal!35!black,title={Before you move on}]
What does \gu{ચાવવું} mean? (\textquotedblleft{}To chew\textquotedblright{}.) Say it once more.
\end{tcolorbox}
\section[\texorpdfstring{\gu{ચૂસવું} (cūsvũ)}{ચૂસવું (cūsvũ)}]{\gu{ચૂસવું} (cūsvũ) — to suck}

\label{lesson:GU-C474-cusvu}



\begin{quote}
Before the new one: say the Gujarati for saliva, then the Gujarati for to chew.
\end{quote}

\subsection*{You'll want to know: \gu{ચૂસવું}}
\textbf{\gu{ચૂસવું}} — \emph{cūsvũ} — \textquotedblleft{}to suck\textquotedblright{}.

\begin{quote}
\gu{બાળક} \gu{અંગૂઠો} \gu{ચૂસે} \gu{છે}.
\end{quote}

\begin{quote}
\emph{bāḷak angūṭho cūse che.}
\end{quote}

\begin{quote}
The baby is sucking its thumb.
\end{quote}

Past, with \gu{તેણે} (he or she …): \gu{તેણે} \gu{કેરી} \gu{ચૂસી}, \emph{teṇe kerī cūsī}.

\begin{cousinweb}
\textbf{\gu{ચૂસવું}} (\emph{cūsvũ}) is inherited from Sanskrit \emph{cūṣati}, \textquotedblleft{}sucks\textquotedblright{}; CDIAL lists the Gujarati verb (CDIAL 4898).
\end{cousinweb}

\begin{grammarlens}[title={What you've built}]
One more word for describing what you lived and felt.
\end{grammarlens}

\subsection*{Your turn}
\begin{itemize}
\raggedright
  \item \emph{Say it:} \emph{cūsvũ}
  \item \emph{Say it:} \emph{cūsvũ}, once more
  \item \emph{Say it:} \emph{bāḷak angūṭho cūse che.} — The baby is sucking its thumb.
  \item \emph{Recall:} say the Gujarati for saliva, then the Gujarati for to chew, then say \emph{cūsvũ} again
\end{itemize}

\begin{tcolorbox}[breakable,colback=teal!4,colframe=teal!35!black,title={Before you move on}]
What does \gu{ચૂસવું} mean? (\textquotedblleft{}To suck\textquotedblright{}.) Say it once more.
\end{tcolorbox}
\section[\texorpdfstring{\gu{થૂંકવું} (thū̃kvũ)}{થૂંકવું (thū̃kvũ)}]{\gu{થૂંકવું} (thū̃kvũ) — to spit}

\label{lesson:GU-C474-thukvu}



\begin{quote}
Before the new one: say the Gujarati for to chew, then the Gujarati for to suck.
\end{quote}

\subsection*{You'll want to know: \gu{થૂંકવું}}
\textbf{\gu{થૂંકવું}} — \emph{thū̃kvũ} — \textquotedblleft{}to spit\textquotedblright{}.

\begin{quote}
\gu{રસ્તા} \gu{પર} \gu{થૂંકવું} \gu{ખરાબ} \gu{ટેવ} \gu{છે}.
\end{quote}

\begin{quote}
\emph{rastā par thū̃kvũ kharāb ṭev che.}
\end{quote}

\begin{quote}
Spitting in the street is a bad habit.
\end{quote}

Past (he …): \gu{તે} \gu{રસ્તા} \gu{પર} \gu{થૂંક્યો}, \emph{te rastā par thū̃kyo}.

\begin{cousinweb}
\textbf{\gu{થૂંકવું}} (\emph{thū̃kvũ}) goes back to a reconstructed \emph{*thukk-}, \textquotedblleft{}spit\textquotedblright{}, an imitative word; CDIAL lists the Gujarati verb (CDIAL 6097).
\end{cousinweb}

\begin{grammarlens}[title={What you've built}]
One more word for describing what you lived and felt.
\end{grammarlens}

\subsection*{Your turn}
\begin{itemize}
\raggedright
  \item \emph{Say it:} \emph{thū̃kvũ}
  \item \emph{Say it:} \emph{thū̃kvũ}, once more
  \item \emph{Say it:} \emph{rastā par thū̃kvũ kharāb ṭev che.} — Spitting in the street is a bad habit.
  \item \emph{Recall:} say the Gujarati for to chew, then the Gujarati for to suck, then say \emph{thū̃kvũ} again
\end{itemize}

\begin{tcolorbox}[breakable,colback=teal!4,colframe=teal!35!black,title={Before you move on}]
What does \gu{થૂંકવું} mean? (\textquotedblleft{}To spit\textquotedblright{}.) Say it once more.
\end{tcolorbox}
\section[\texorpdfstring{\gu{તાળવું} (tāḷvũ)}{તાળવું (tāḷvũ)}]{\gu{તાળવું} (tāḷvũ) — the palate}

\label{lesson:GU-C474-talvu}



\begin{quote}
Before the new one: say the Gujarati for to suck, then the Gujarati for to spit.
\end{quote}

\subsection*{You'll want to know: \gu{તાળવું}}
\textbf{\gu{તાળવું}} — \emph{tāḷvũ} — \textquotedblleft{}the palate\textquotedblright{}.

\begin{quote}
\gu{ગરમ} \gu{ચાથી} \gu{મારું} \gu{તાળવું} \gu{બળી} \gu{ગયું}.
\end{quote}

\begin{quote}
\emph{garam cāthī mārũ tāḷvũ baḷī gayũ.}
\end{quote}

\begin{quote}
The hot tea burned the roof of my mouth.
\end{quote}

A neuter noun.

\begin{cousinweb}
\textbf{\gu{તાળવું}} (\emph{tāḷvũ}) is inherited from Sanskrit \emph{tālu}, \textquotedblleft{}palate\textquotedblright{}; CDIAL lists the Gujarati word, \textquotedblleft{}palate, crown of head\textquotedblright{} (CDIAL 5803).
\end{cousinweb}

\begin{grammarlens}[title={What you've built}]
One more word for describing what you lived and felt.
\end{grammarlens}

\subsection*{Your turn}
\begin{itemize}
\raggedright
  \item \emph{Say it:} \emph{tāḷvũ}
  \item \emph{Say it:} \emph{tāḷvũ}, once more
  \item \emph{Say it:} \emph{garam cāthī mārũ tāḷvũ baḷī gayũ.} — The hot tea burned the roof of my mouth.
  \item \emph{Recall:} say the Gujarati for to suck, then the Gujarati for to spit, then say \emph{tāḷvũ} again
\end{itemize}

\begin{tcolorbox}[breakable,colback=teal!4,colframe=teal!35!black,title={Before you move on}]
What does \gu{તાળવું} mean? (\textquotedblleft{}The palate\textquotedblright{}.) Say it once more.
\end{tcolorbox}
